\section{Scope and open theorem program}
\label{sec:scope}

\subsection{Scope of the present results}

The paper establishes a conditional principle and a model-level audit. Its scope is defined by the following boundaries:
\begin{enumerate}
\item The translational inverse is global for \(\gamma\neq0\) and velocity-independent force; the full inverse additionally requires a unique phase branch.
\item The certificate \(q(\bm r^+)<1\) proves uniqueness for an individual phase update. A nontrivial invariant domain on which it holds for all time remains to be characterized.
\item Jacobian and differential-entropy identities apply on smooth strata. The finite anchor cutoffs are piecewise smooth.
\item The phrase readout is a deterministic geometric clustering with an auxiliary score. It is not required to be a complete code.
\item The round-trip experiment measures end-to-end float64 recovery, not an independently estimated condition number.
\item The model separates exact-state, observational, and numerical information. A thermodynamic theory would additionally require energy, reservoirs, temperature, and a physical measure.
\item The equations define a mathematical system. Applying the principle to a physical domain requires an independently justified complete state, calibrated parameters, and prospective tests.
\end{enumerate}
These boundaries delimit the theorem without reducing its generality as an audit method.

\subsection{Global phase domain}

The phase Jacobian is an identity minus a signed weighted graph Laplacian. The principal open theorem is to characterize the largest domain on which \(G_{\bm r}\) is globally bijective and invariant under the full advance. Promising routes include:
\begin{enumerate}
\item spectral bounds sharper than the \(\ell_\infty\) certificate \(q<1\);
\item degree-theoretic continuation of the torus map;
\item explicit branch-bifurcation analysis beyond the two-clock reduction;
\item invariant minimum-separation estimates for the coupled translational dynamics.
\end{enumerate}

\subsection{Invertible boundaries}

The clamp in \cref{prop:wall} can be replaced by a distinction-preserving boundary:
\begin{itemize}
\item periodic position on a two-torus;
\item exact specular reflection retaining overshoot;
\item an enlarged state that records the boundary impulse and traversed distance.
\end{itemize}
Each option yields a different but explicitly auditable state space.

\subsection{Continuous-time realization}

A smooth continuous counterpart is
\begin{subequations}
\begin{align}
\dot{\bm r}_i&=\bm v_i,\\
\dot{\bm v}_i&=-\kappa\bm v_i+m_i^{-1}\bm F_i(\bm r,\bm\theta,\bm\phi),\\
\dot\phi_i&=\omega_i+\beta\sum_{j\neq i}
\frac{\sin(\phi_j-\phi_i)}
{\norm{\bm r_i-\bm r_j}+\varepsilon}.
\end{align}
\label{eq:continuous}
\end{subequations}
For a globally Lipschitz vector field, uniqueness of solutions gives an invertible finite-time flow wherever solutions exist. Liouville's general formula,
\[
\det D\Phi_t(x_0)
=
\exp\left(
\int_0^t\nabla\cdot f(x_s)\,\dd s
\right),
\]
then relates volume contraction to the divergence without identifying trajectories. A smooth taper at the anchor cutoff would place the complete realization inside this theorem.

\subsection{Recovery bounds}

The computed round-trip curves motivate a trajectory-level error theorem. Let
\[
\kappa_T(x_t):=\norm{D\Psi^{-T}(x_t)}.
\]
A useful bound would expose the velocity multiplier, phase singular values, and force Lipschitz constants:
\begin{equation}
\norm{\widehat x_{t-T}-x_{t-T}}
\le
C_T\epsilon,
\qquad
C_T
\lesssim
\prod_{s=0}^{T-1}
\frac{c_s}
{\abs{\gamma}\,\sigma_{\min}(D_\phi G_s)}.
\label{eq:horizon-bound}
\end{equation}
Given a tolerance \(\tau\), the largest \(T\) for which \(C_T\epsilon\le\tau\) would provide a certified numerical access horizon.

\subsection{Observability from histories}

The state-sufficiency result invites a delay-coordinate analysis. For a readout \(y_t=h(x_t)\), determine when
\[
(y_t,y_{t-\tau},\ldots,y_{t-(m-1)\tau})
\]
embeds the reachable set, how large \(m\) must be, and how inverse conditioning scales with noise. This would connect exact-state theory to realistic partial instrumentation.

\subsection{Description-level dynamics}

The closure criterion \cref{prop:closure} should be evaluated directly for the phrase readout and for task-specific alternatives. When closure fails, one can quantify the memory order required for a predictive description. When it holds in both directions, the induced macro-law inherits invertibility even though the readout is many-to-one.

\subsection{Research sequence}

\[
\boxed{
\text{phase domain}
\longrightarrow
\text{invertible boundary}
\longrightarrow
\text{recovery bound}
\longrightarrow
\text{history observability}
\longrightarrow
\text{domain tests}.
}
\]
This sequence advances the principle from exact one-step structure to certified long-horizon and domain-specific use.
